% =============================================================
% §5 Evaluation
% =============================================================
\section{Evaluation}
\label{sec:eval}

\subsection{Experimental Setup}
\label{sec:eval-setup}

\textbf{Simulation platform.} All experiments use the
simulator of \S\ref{sec:simulator}, which implements the five
AIGC physics as milestones M1--M4.

\textbf{Hardware configuration.} One cloud-class server
(200~TFLOPS, 128~GB VRAM, A100-80GB-equivalent) plus 1--7
heterogeneous edge nodes (10--50~TFLOPS, 16--64~GB VRAM,
modeled after Jetson AGX, T4, and edge A100 deployments).
Network latency 30--65~ms for cloud--edge links, 8--23~ms for
edge--edge links.

\textbf{Workload.} We use a log-normal request length
distribution matching the published statistics of the Azure LLM
Inference Trace 2023~\cite{azure2023llmtrace}: prompt length
$\sim \text{LogNormal}(\mu{=}5.5, \sigma{=}1.0)$, output
length $\sim \text{LogNormal}(\mu{=}4.5, \sigma{=}1.2)$.
Request arrival follows a Poisson process at rate $\lambda$.
Each request is assigned uniformly to one of three models
(LLaMA-7B, LLaMA-13B, LLaMA-70B-INT8) with parameters
calibrated to vLLM and TensorRT-LLM benchmark numbers.

\textbf{Baselines (11 schedulers).} We compare against a broad
set:
\begin{itemize}
\item \textbf{RoundRobin (RR)}---load-agnostic baseline.
\item \textbf{LeastLoaded (LL)}---load-aware greedy.
\item \textbf{ShortestQueue (SQ)}---queue-aware greedy.
\item \textbf{HEFT}~\cite{topcuoglu2002heft}---the standard
  DAG-scheduling baseline.
\item \textbf{GA}~\cite{holland1975adaptation},
  \textbf{PSO}~\cite{kennedy1995pso}---metaheuristic search.
\item \textbf{NSGA-II}~\cite{yu2025llmrouting}---multi-objective
  genetic search (non-dominated sorting + crowding distance)
  over per-window assignments, with two objectives matching
  our axes (completion-time proxy, energy estimate); the
  knee-point solution of the final Pareto front is executed.
\item \textbf{A3C-R2N2}~\cite{tuli2022a3cr2n}---RL baseline
  using asynchronous advantage actor-critic with a
  GRU+residual encoder but with \textbf{generic (non-AIGC)
  state and reward}.
\item \textbf{GNN}---a graph-neural-network variant of our
  scheduler that encodes AIGC physics as \textbf{edge features}
  in a task--server bipartite graph, similar to
  Decima~\cite{mao2019decima}.
\item \textbf{CS-UCB}~\cite{yang2024perllm}---a PerLLM-style
  constraint-satisfaction UCB bandit, the closest prior work
  (\S\ref{sec:related}): per-request server selection by
  maximum UCB among feasible arms, with reward
  $-E_{\text{est}} + \lambda f(y)$ (negative energy estimate
  plus normalized constraint slack).
\item \textbf{Ours (RL)}---our AIGC-aware PPO scheduler with
  explicit AIGC state features and reward shaping.
\end{itemize}

All eleven schedulers are implemented natively in Python
against
the simulator's uniform \texttt{scheduler.schedule()}
interface; no external scheduler library is reused. The
NSGA-II baseline mirrors the multi-objective cloud-edge
routing setting of~\cite{yu2025llmrouting}. HEFT
follows the original upward-rank
formulation~\cite{topcuoglu2002heft}; A3C-R2N2 reimplements
the GRU + residual-skip network of~\cite{tuli2022a3cr2n}; the
GNN variant follows Decima-style graph
encoding~\cite{mao2019decima}; CS-UCB reimplements PerLLM's
constraint-satisfaction UCB~\cite{yang2024perllm} at our
per-task granularity, with per-(model, phase, server) arm
statistics mirroring its per-service personalization.
Table~\ref{tab:hyperparams}
lists each scheduler's hyperparameters and search/training
budget; all implementations ship with our open-source release.

\begin{table}[t]
  \centering
  \caption{Scheduler implementations and hyperparameters.
    RR/LL/SQ = RoundRobin, LeastLoaded, ShortestQueue.}
  \label{tab:hyperparams}
  \small
  \setlength{\tabcolsep}{4pt}
  \begin{tabular}{@{}l@{\hspace{6pt}}p{0.64\columnwidth}@{}}
    \toprule
    Scheduler & Configuration \\
    \midrule
    RR / LL / SQ & parameter-free greedy rules \\
    HEFT & upward rank + earliest finish time; parameter-free \\
    GA & population 30, 50 generations per scheduling window;
      crossover 0.8, mutation 0.1, elitism 0.1 \\
    PSO & 30 particles, 30 iterations per window; inertia
      $w{=}0.7$, $c_1{=}c_2{=}1.5$ \\
    NSGA-II & population 30, 30 generations per window;
      crossover 0.8, mutation 0.1; executes the knee point of
      the final non-dominated front \\
    A3C-R2N2 & GRU cell + residual skip (hidden 128); Adam
      $3{\times}10^{-4}$, $\gamma{=}0.95$, entropy 0.05, value
      coef.\ 0.5; update every 32 transitions (single epoch,
      no clipping, per~\cite{tuli2022a3cr2n}) \\
    GNN & 1-hop attention message passing on the task--server
      bipartite graph (hidden 64); same PPO recipe and budget
      as ours \\
    CS-UCB & arms = (model, phase, server); feasibility filter
      via admission check; exploration $\delta{=}0.5$,
      constraint weight $\lambda{=}1.0$ (grid-tuned over
      $\delta \in \{0.1,\dots,2\}$, $\lambda \in \{0.5,1\}$) \\
    RL (ours) & MLP $256{\to}128{\to}64$; PPO clip
      $\epsilon{=}0.2$, GAE $\gamma{=}0.95$,
      $\lambda{=}0.90$; Adam $3{\times}10^{-4}$; entropy
      $0.05{\to}0.01$ after pretrain; update every 32
      transitions $\times$ 4 epochs \\
    \addlinespace
    \multicolumn{2}{@{}p{0.97\columnwidth}@{}}{\emph{Training
      budget (all learned schedulers):} 20 pretraining
      episodes + 30 random warm-up dispatches; online updates
      continue during evaluation.} \\
    \bottomrule
  \end{tabular}
\end{table}

\textbf{Metrics.} Following standard LLM serving benchmarks:
\slo{} attainment (fraction meeting TTFT $\le$ 2.0~s and TPOT
$\le$ 100~ms/token); TTFT and TPOT percentiles;
goodput; total energy; and energy per
token (J/tok), using the calibrated power model of
\S\ref{sec:simulator}.

\textbf{Statistical method.} All experiments use 30
independent seeds; we report mean~$\pm$~std and test
significance with two-sided Mann-Whitney U (simulation
outputs are non-Gaussian). Thresholds:
$*$ $p<0.05$, $**$ $p<0.01$, $***$ $p<0.001$.

\subsection{Main Results: Pareto Dominance}
\label{sec:main-results}

We evaluate all eleven schedulers on the canonical configuration:
5 servers (1 cloud + 4 heterogeneous edges), Poisson arrival
$\lambda{=}2$ req/s, 100 task units ($\approx$50 inference
requests). Results are averaged over 30 runs.

\textbf{Headline finding.} Our AIGC-aware RL scheduler
simultaneously achieves the highest \slo{} attainment
\textbf{and} the lowest energy per token, \textbf{Pareto-dominating
all ten baselines on mean performance} in the (\slo, energy
efficiency) plane
(Table~\ref{tab:main}, Fig.~\ref{fig:pareto}).

\begin{table}[t]
  \centering
  \caption{Main comparison at edge${=}5$, $\lambda{=}2$ req/s,
    $N{=}30$. Mksp.\ = makespan (s); E = total energy (kJ);
    \etok{} = energy per token (J); TTFT P95 in seconds.
    Bold = column-best.}
  \label{tab:main}
  \small
  \setlength{\tabcolsep}{4pt}
  \begin{tabular}{lrrrrr}
    \toprule
    Scheduler & Mksp.(s) & \slo$\uparrow$ & E(kJ)$\downarrow$
              & \etok(J)$\downarrow$ & TTFT P95 \\
    \midrule
    RoundRobin    & 306.2 & 0.163 & 64.0 & 7.29 & 174.6 \\
    LeastLoaded   & 125.5 & 0.239 & 26.8 & 2.94 & 30.8 \\
    ShortestQueue & 123.8 & 0.234 & 26.9 & 2.96 & 25.4 \\
    HEFT          & 126.7 & 0.169 & 25.9 & 2.83 & 35.7 \\
    GA            & 132.4 & 0.225 & 26.9 & 2.94 & 30.5 \\
    PSO           & 129.6 & 0.235 & 25.6 & 2.81 & 41.0 \\
    A3C-R2N2      & \textbf{120.6} & 0.214 & 26.7 & 2.93 & 28.8 \\
    GNN           & 125.5 & 0.210 & 26.7 & 2.93 & 32.2 \\
    CS-UCB        & 129.9 & 0.193 & 26.0 & 2.84 & 37.3 \\
    NSGA-II       & 126.2 & 0.289 & 25.5 & 2.79 & 32.0 \\
    \textbf{RL (ours)} & 122.2 & \textbf{0.302}
                       & \textbf{23.7} & \textbf{2.61}
                       & \textbf{25.2} \\
    \bottomrule
  \end{tabular}
\end{table}

The proposed RL
scheduler is best on every metric except makespan, where
A3C-R2N2 completes marginally earlier (120.6~s vs.\ 122.2~s, a
1.3\% gap); makespan is not among our objectives
(\S\ref{sec:sched-problem}) and the gap does not affect the
Pareto relation on the (\slo, \etok) plane.

\begin{figure}[t]
  \centering
  \includegraphics[width=\columnwidth]{fig4_pareto.pdf}
  \caption{Pareto plot of \slo{} attainment vs.\ energy per
    token at the canonical configuration (edge${=}5$,
    $\lambda{=}2$ req/s, $N{=}30$). The ``ideal'' corner is
    upper-left. Our RL scheduler sits alone in this region.
    Single-objective baselines cluster to the lower-right
    with markedly lower \slo; the closest challenger,
    NSGA-II, trails on both axes.
    RoundRobin is an extreme outlier off-chart.}
  \label{fig:pareto}
\end{figure}

\textbf{Statistical significance.} Table~\ref{tab:significance}
reports Mann-Whitney U tests
comparing our RL scheduler against each of the six strongest
baselines.

\begin{table}[t]
  \centering
  \caption{$p$-values vs.\ six strongest baselines (two-sided
    Mann-Whitney U, $N{=}30$). Bold = $p<0.05$.}
  \label{tab:significance}
  \scriptsize
  \setlength{\tabcolsep}{3pt}
  \begin{tabular}{lcccccc}
    \toprule
    Metric & PSO & SQ & A3C-R2N2 & GNN & CS-UCB & NSGA-II \\
    \midrule
    \slo{} attain.   & \textbf{0.045}$^*$ & \textbf{0.046}$^*$
                     & \textbf{0.006}$^{**}$ & \textbf{0.006}$^{**}$
                     & \textbf{0.002}$^{**}$ & 0.64 \\
    Total energy     & 0.36 & 0.07 & 0.09 & 0.10 & 0.21 & 0.27 \\
    Energy/token     & 0.11 & \textbf{0.010}$^{**}$
                     & \textbf{0.028}$^*$ & \textbf{0.016}$^*$
                     & 0.066 & 0.12 \\
    \bottomrule
  \end{tabular}
\end{table}

\slo{} attainment is significantly higher in our scheduler
against all five single-objective baselines ($p<0.05$, with
$p<0.01$ for both
DRL baselines and the bandit). Energy efficiency is
significantly better
against three of six (PSO, $p=0.11$, and CS-UCB, $p=0.066$,
are directionally consistent but
do not reach $p<0.05$). The one baseline our scheduler does
not separate from statistically at this operating point is
NSGA-II ($p=0.64$ \slo, $p=0.12$ \etok): our scheduler is
better on both means, but the gaps are within noise at
$N{=}30$; we return to this comparison---including its
sweet-spot significance and cost asymmetry---below. Total
energy shows a consistent
$\le 10\%$ advantage but does not reach individual
significance, reflecting the expected correlation between
makespan and energy.

\textbf{Why does AIGC-aware RL achieve lower energy?}
Counterintuitively, our scheduler---with access to high-power
cloud resources---uses \emph{less} energy than baselines that
default to edge-heavy behavior. We attribute this to three
learned behaviors: (i)~\emph{shorter makespan amplifies idle
savings} (cluster idle power is $\approx$150~W);
(ii)~\emph{cold-load avoidance}---co-locating same-model
requests eliminates loads that keep GPUs active without
producing tokens; (iii)~\emph{sibling affinity}---decode
placed beside its prefill avoids the active-but-idle
KV-transfer interval. \textbf{With energy as a co-objective,
these behaviors make the AIGC-aware design Pareto-dominant.}

\textbf{Comparison to A3C-R2N2 (controlled RL backbone).}
A3C-R2N2 uses the same actor-critic framework but with (a) a
GRU+residual encoder rather than our MLP, and (b) generic
state and reward without AIGC components. This is the
closest-controlled comparison. A3C-R2N2 achieves \slo${} =
0.214 \pm 0.07$ and \etok${} = 2.93$~J/tok, compared to our
$0.302 \pm 0.07$ and $2.61$~J/tok. The differences are
significant at $p<0.01$ (\slo) and $p<0.05$ (\etok). Since the
backbones are equivalent (both actor-critic with adequate
pretraining), this $41\%$ \slo{} improvement and $11\%$
energy efficiency improvement is \textbf{directly attributable
to AIGC-aware state and reward design}.

\textbf{Comparison to GNN (controlled state encoding).} The
GNN variant encodes AIGC physics as \textbf{edge features} in
a task--server bipartite graph, an alternative to our MLP
with flat AIGC features. Despite the architectural
sophistication, GNN achieves \slo${} = 0.210 \pm 0.08$ and
\etok${} = 2.93$~J/tok, essentially matching A3C-R2N2 and
significantly worse than our scheduler ($p<0.01$ \slo,
$p<0.05$ \etok). We attribute this gap to \textbf{training
data efficiency}: the GNN's relational reasoning capacity
requires more samples to specialize than the pretrain budget
(20 episodes) provides.

\textbf{Comparison to CS-UCB (PerLLM-style bandit).} CS-UCB
shares our per-request granularity and energy objective, and
its constraint-satisfaction mechanism mirrors
PerLLM~\cite{yang2024perllm}. Yet it lands in the
baseline cluster (\slo{} 0.193, \etok{} 2.84): significantly
below our scheduler on \slo{} ($p{=}0.002$), directionally
worse on energy ($p{=}0.066$). The cause is structural:
stationary per-arm means cannot track
\emph{state-dependent} costs---model residency and batch
slots depend on the cluster's evolving
state, which a context-free bandit averages away while a deep
policy conditions on it. Constraint filtering does not
substitute for state-conditioned learning.

\textbf{Comparison to NSGA-II (multi-objective search).}
NSGA-II is by far the strongest baseline and deserves a candid
account. Searching per-window assignments against the same two
objectives we optimize, it reaches \slo{} 0.289 and \etok{}
2.79 at the canonical point---our scheduler is better on both
means (0.302, 2.61) but neither gap is individually
significant at $N{=}30$ ($p=0.64$, $p=0.12$). Three
observations complete the picture. First, in the
\textbf{moderate-load sweet spot} the separation is real: at
$\lambda{=}1$ our scheduler reaches \slo{} 0.447 vs.\
NSGA-II's 0.352 ($+27\%$ relative, $p=0.009$). Second, our
scheduler retains the \textbf{lower energy per token in every
one of the eight configurations} we ran NSGA-II on, with
margins growing from parity at $\lambda{=}0.5$ to 7--9\%
under load. Third, the comparison is asymmetric in
\textbf{scheduling-time cost}: NSGA-II re-runs a population
search at dispatch time ($5.4\times$ our total wall-clock;
$\approx$70~ms of search per dispatched task, measured),
whereas our policy decides by a sub-millisecond forward pass
after a one-off pretrain---in latency-sensitive serving,
dispatch-time search inflates the very TTFT it is
protecting. We therefore read NSGA-II as the strongest
\emph{offline} competitor; its strength at saturation
(\S\ref{sec:load}) sharpens the regime characterization
(C4).

\subsection{Topology Sensitivity}
\label{sec:topology}

We vary the number of edge servers from 3 to 7 (always with
one cloud) at fixed $\lambda{=}2$ req/s and uniform model mix.
Table~\ref{tab:topology} reports \slo{} and \etok{} for the
four strongest baselines and our scheduler.

\begin{table}[t]
  \centering
  \caption{\slo{} and \etok{} (J/tok) vs.\ edge count.
    $\bigstar$ = column-best across all 11 schedulers.}
  \label{tab:topology}
  \scriptsize
  \setlength{\tabcolsep}{2.5pt}
  \begin{tabular}{lrrrrrrr}
    \toprule
    Config (1 cloud +)
      & PSO & SQ & A3C-R2N2 & GNN & CS-UCB & NSGA-II & \textbf{RL} \\
    \midrule
    3 edges --- \slo
      & 0.273 & 0.277 & 0.246 & 0.241 & 0.243 & 0.296
      & \textbf{0.314}$^\bigstar$ \\
    3 edges --- \etok
      & 2.19 & 2.32 & 2.28 & 2.24 & 2.24 & 2.16
      & \textbf{1.99}$^\bigstar$ \\
    5 edges --- \slo
      & 0.235 & 0.234 & 0.214 & 0.210 & 0.193 & 0.289
      & \textbf{0.302}$^\bigstar$ \\
    5 edges --- \etok
      & 2.81 & 2.96 & 2.93 & 2.93 & 2.84 & 2.79
      & \textbf{2.61}$^\bigstar$ \\
    7 edges --- \slo
      & 0.241 & 0.218 & 0.193 & 0.250 & 0.152 & 0.277
      & \textbf{0.282}$^\bigstar$ \\
    7 edges --- \etok
      & 3.25 & 3.34 & 3.35 & 3.35 & 3.35 & 3.19
      & \textbf{3.00}$^\bigstar$ \\
    \bottomrule
  \end{tabular}
\end{table}

Across all three topology configurations, our RL scheduler is
the best on every metric. Against the single-objective
baselines, RL leads \slo{} by 13--28\% relative and energy by
7--9\%. The strongest challenger, NSGA-II, narrows the
\slo{} gap to 2--6\% relative but pays 6--8\% more energy per
token at every topology---the same asymmetry we observe at
the canonical point.

\textbf{Energy advantage is more \emph{configuration-stable}
than \slo.} The relative energy improvement sits in a narrow
band [7\%, 9\%] across all topology configurations, while the
\slo{} improvement varies more (13--28\%). This suggests
\textbf{energy efficiency captures the ``intrinsic'' benefit
of AIGC-aware scheduling more cleanly than \slo}, since \slo{}
is sensitive to cluster contention dynamics whereas energy
directly reflects per-task efficiency.

\subsection{Load Sensitivity}
\label{sec:load}

We sweep the Poisson arrival rate
$\lambda \in \{0.5, 1, 2, 4, 8\}$ req/s at fixed edge${=}5$ to
characterize how scheduler behavior changes from
under-utilization to over-saturation.
Table~\ref{tab:slo-vs-lambda}
reports \slo{} attainment and energy per token across the
sweep.

\begin{table}[t]
  \centering
  \caption{Load sweep: \slo{} attainment (top) and energy per
    token (J/tok, bottom) across arrival rates. $\bigstar$
    = strict winner; $\dagger$ = tied with PSO.}
  \label{tab:slo-vs-lambda}
  \scriptsize
  \setlength{\tabcolsep}{2.5pt}
  \begin{tabular}{rrrrrrrl}
    \toprule
    $\lambda$ & PSO & SQ & A3C-R2N2 & GNN & CS-UCB & NSGA-II
      & \textbf{RL} \\
    \midrule
    \multicolumn{8}{@{}l}{\emph{\slo{} attainment}} \\
    0.5 & 0.596 & 0.339 & 0.359 & 0.500 & 0.462 & 0.577
        & \textbf{0.596}$^\dagger$ \\
    1.0 & 0.297 & 0.265 & 0.227 & 0.237 & 0.214 & 0.352
        & \textbf{0.447}$^\bigstar$ \\
    2.0 & 0.235 & 0.234 & 0.214 & 0.210 & 0.193 & 0.289
        & \textbf{0.302}$^\bigstar$ \\
    4.0 & 0.253 & 0.205 & 0.191 & 0.207 & 0.215
        & \textbf{0.267} & 0.214 \\
    8.0 & 0.216 & 0.167 & 0.156 & 0.193 & 0.187
        & \textbf{0.219} & 0.199 \\
    \addlinespace
    \multicolumn{8}{@{}l}{\emph{Energy per token (J/tok)}} \\
    0.5 & 3.00 & 3.49 & 3.38 & 3.13 & 3.19 & 2.97
        & \textbf{2.96}$^\bigstar$ \\
    1.0 & 2.77 & 3.00 & 3.07 & 2.95 & 2.93 & 2.84
        & \textbf{2.64}$^\bigstar$ \\
    2.0 & 2.81 & 2.96 & 2.93 & 2.93 & 2.84 & 2.79
        & \textbf{2.61}$^\bigstar$ \\
    4.0 & 2.85 & 2.97 & 2.94 & 2.92 & 2.86 & 2.78
        & \textbf{2.58}$^\bigstar$ \\
    8.0 & 2.94 & 2.89 & 2.98 & 2.91 & 2.99 & 2.78
        & \textbf{2.52}$^\bigstar$ \\
    \bottomrule
  \end{tabular}
\end{table}

\textbf{Key observations.}
(1)~\emph{Energy efficiency is RL-dominant across all loads.}
RL achieves the best energy per token at every $\lambda$
(5/5), with a margin over the next-best baseline that grows
monotonically with load---from parity at $\lambda{=}0.5$ to
$9\%$ at $\lambda{=}8$.
(2)~\emph{\slo{} dominance is concentrated in the
moderate-load regime} ($\lambda \in [0.5, 2]$), where RL is
strict winner or tied at every point; at $\lambda{=}1$ the
improvement is $27\%$ relative over the strongest baseline
(NSGA-II, $p=0.009$) and $50\%$ over PSO.
(3)~\emph{At very high load ($\lambda \ge 4$)}, the cluster
approaches its physical throughput ceiling and offline
multi-objective search takes over the \slo{} axis: NSGA-II is
the strict winner at $\lambda{=}4$ and $\lambda{=}8$, with
PSO close behind---batched search exploits the homogenized
queue state better than any incremental policy---and even
CS-UCB reaches
RL's \slo{} level at $\lambda{=}4$ (0.215 vs.\ 0.214).
\textbf{However, RL's
energy advantage persists and peaks} in this regime
(Table~\ref{tab:slo-vs-lambda}, bottom block).

\textbf{When does AIGC-aware design matter most?} We
characterize three regimes:
\textbf{under-loaded} ($\lambda \le 1$, spare
capacity)---best or tied on both axes;
\textbf{sweet spot} ($1 \le \lambda \le 2$)---strict Pareto
winner, with the $\lambda{=}1$ \slo{} lead significant even
against NSGA-II;
\textbf{saturated} ($\lambda \ge 4$, throughput
ceiling)---energy-only edge, with the \slo{} axis ceded to
offline multi-objective search. This nuanced finding is more
defensible than a flat ``we
win everywhere'' claim. The cleanest statement is:
\emph{AIGC-aware RL achieves the lowest energy per token at
every load---with a margin that grows with load---and strict
\slo{}+energy
Pareto dominance in the moderate-load operating regime that
characterizes most production inference deployments.}

\subsection{Workload Composition Sensitivity}
\label{sec:workload}

We test two model mixes at edge${=}5$, $\lambda{=}2$: \textbf{uniform}
(1/3 each of LLaMA-7B/13B/70B) and \textbf{large-heavy} (70\%
LLaMA-70B-INT8, biased toward memory-bound requests).
Table~\ref{tab:workload} reports the results.

\begin{table}[t]
  \centering
  \caption{Pareto performance by workload mix.}
  \label{tab:workload}
  \small
  \begin{tabular}{llrr}
    \toprule
    Mix & Scheduler & \slo{} & \etok{} \\
    \midrule
    uniform & PSO            & 0.235 & 2.81 \\
    uniform & ShortestQueue  & 0.234 & 2.96 \\
    uniform & A3C-R2N2       & 0.214 & 2.93 \\
    uniform & GNN            & 0.210 & 2.93 \\
    uniform & CS-UCB         & 0.193 & 2.84 \\
    uniform & NSGA-II        & 0.289 & 2.79 \\
    uniform & \textbf{RL (ours)}
            & \textbf{0.302}$^\bigstar$ & \textbf{2.61}$^\bigstar$ \\
    \midrule
    large-heavy & PSO            & 0.093 & 3.84 \\
    large-heavy & ShortestQueue  & 0.105 & 3.84 \\
    large-heavy & A3C-R2N2       & 0.107 & 3.79 \\
    large-heavy & GNN            & 0.091 & 3.79 \\
    large-heavy & CS-UCB         & 0.068 & 3.84 \\
    large-heavy & NSGA-II        & \textbf{0.122} & 3.68 \\
    large-heavy & \textbf{RL (ours)}
                & 0.098 & \textbf{3.58}$^\bigstar$ \\
    \bottomrule
  \end{tabular}
\end{table}

In the \textbf{uniform mix}, our RL is the strict Pareto winner.
In the \textbf{large-heavy mix}, the picture shifts: all
schedulers collapse in \slo{} to a low band (0.07--0.12,
vs.\ 0.19--0.30 under uniform)
because LLaMA-70B-INT8 (70~GB weights) physically fits
\emph{only on the cloud server}. The cloud queue saturates
regardless of scheduler choice; within the band, NSGA-II's
offline search attains the best \slo{} (0.122), echoing its
saturation advantage in \S\ref{sec:load}. \textbf{Energy
efficiency is
the only broadly differentiable axis}, where RL still leads
(3.58 vs.\ NSGA-II's 3.68 and 3.79--3.84 for the rest).
\emph{When the
workload exceeds cluster capacity at a specific tier,
AIGC-aware scheduling degrades to pure energy optimization}
(discussed in \S\ref{sec:discuss}).

\subsection{Component Ablation}
\label{sec:ablation}

To isolate which of our RL scheduler's components contribute
to its Pareto dominance, we systematically ablate 12 components
by disabling one at a time, holding all else equal. All
ablations use edge${=}5$, $\lambda{=}2$ req/s, uniform mix,
$N{=}30$ runs. Table~\ref{tab:ablation} reports all twelve
configurations.

\textbf{Ablation taxonomy.} We group ablations into four
categories: \emph{Physical layer} (\texttt{no\_batching});
\emph{PPO internals} (\texttt{no\_action\_mask},
\texttt{no\_gae}, \texttt{no\_pretrain}, \texttt{no\_entropy});
\emph{AIGC reward} (\texttt{no\_warm\_reward},
\texttt{no\_batch\_reward}, \texttt{no\_affinity\_reward},
\texttt{no\_cloud\_overuse},
\texttt{no\_all\_aigc\_rewards}); and \emph{AIGC state}
(\texttt{no\_aigc\_state}, \texttt{no\_aigc\_full}).

\begin{table*}[t]
  \centering
  \caption{Component ablation. $\Delta$ shown relative to Full
    RL ($N{=}30$, edge${=}5$, $\lambda{=}2$, uniform mix).
    $p$-values from Mann-Whitney U vs.\ Full RL on \slo.
    Bold = $p<0.05$.}
  \label{tab:ablation}
  \small
  \begin{tabular}{llrrrrr}
    \toprule
    Category & Configuration & \slo{} & $\Delta$\slo\%
             & \etok{} & $\Delta$\etok\% & $p$(\slo) \\
    \midrule
    --- & \textbf{Full RL (ours)}
        & \textbf{0.302} & ---
        & \textbf{2.61} & --- & --- \\
    \midrule
    \textbf{Physical}
      & \textbf{$-$ Continuous batching}
      & \textbf{0.051} & \textbf{$-$83.2\%}
      & \textbf{5.75} & \textbf{$+$120.5\%}
      & \textbf{$<$0.001}$^{***}$ \\
    \textbf{PPO}
      & \textbf{$-$ Pretraining (20$\to$0 ep)}
      & \textbf{0.186} & \textbf{$-$38.4\%}
      & \textbf{2.96} & \textbf{$+$13.5\%}
      & \textbf{$<$0.001}$^{***}$ \\
    PPO & $-$ Action masking
      & 0.274 & $-$9.3\% & 2.66 & $+$2.1\% & 0.39 \\
    PPO & $-$ GAE (use Monte-Carlo return)
      & 0.292 & $-$3.3\% & 2.61 & $-$0.1\% & 0.65 \\
    PPO & $-$ Entropy regularization
      & 0.305 & $+$0.9\% & 2.54 & $-$2.8\% & 0.92 \\
    AIGC reward & $-$ Warm bonus
      & 0.300 & $-$0.7\% & 2.60 & $-$0.4\% & 0.87 \\
    AIGC reward & $-$ Batch bonus
      & 0.311 & $+$3.1\% & 2.59 & $-$0.8\% & 0.87 \\
    AIGC reward & $-$ Affinity bonus
      & 0.287 & $-$4.9\% & 2.64 & $+$1.3\% & 0.57 \\
    AIGC reward & $-$ Cloud-overuse penalty
      & 0.309 & $+$2.2\% & 2.60 & $-$0.5\% & 0.96 \\
    AIGC reward (joint) & $-$ All AIGC rewards
      & 0.302 & $+$0.0\% & 2.62 & $+$0.5\% & 0.95 \\
    AIGC state & $-$ AIGC state features
      & 0.295 & $-$2.2\% & 2.55 & $-$2.1\% & 0.81 \\
    AIGC (joint) & $-$ All AIGC design
      & 0.306 & $+$1.3\% & 2.52 & $-$3.3\% & 0.99 \\
    \bottomrule
  \end{tabular}
\end{table*}

\textbf{Findings: two dominant components and a striking null
result.}

\textbf{Finding 1: Continuous batching is the physical foundation}
($p<0.001$, $\Delta$\slo${} = -83.2\%$, $\Delta$\etok${} =
+120.5\%$). Removing batching collapses \slo{} to 5\% and
more than doubles energy per token, validating the emphasis
of Orca~\cite{yu2022orca} and
Sarathi-Serve~\cite{agrawal2024sarathi} on iteration-level
batching as foundational infrastructure.

\textbf{Finding 2: Adequate pretraining is essential}
($p<0.001$, $\Delta$\slo${} = -38.4\%$, $\Delta$\etok${} =
+13.5\%$). Without pretrain, placement collapses toward the
cloud server---aligning with observations
(Decima~\cite{mao2019decima},
RLScheduler~\cite{zhang2020rlscheduler}) that warm-starting
is essential in stochastic scheduling environments.

\textbf{Finding 3: Individual AIGC components show no
significant effect} ($p > 0.39$ for all). Disabling any single
AIGC component yields within-noise changes
($\pm 5\%$)---consistent with PPO, given adequate base
features and pretraining, \emph{implicitly learning}
AIGC-aware behavior, explicit shaping providing redundant
signal.

\textbf{Finding 4 (most striking): Removing \emph{all} AIGC
design also has no effect} ($p=0.99$ for \slo,
$\Delta$\etok${} = -3.3\%$). The combined ablation
\texttt{no\_aigc\_full} shows essentially identical \slo{}
performance to Full RL, and \emph{slightly better} energy
efficiency. Combined with \S\ref{sec:main-results}'s finding
that our scheduler Pareto-dominates ten baselines---yet its
AIGC-specific design components are individually and jointly
ablatable---this paints a more subtle picture:
\textbf{the contribution of ``AIGC awareness'' is captured in
the training procedure + simulator physics + RL backbone
choice, rather than in any single reward or state component}
(Fig.~\ref{fig:ablation}).

\begin{figure*}[t]
  \centering
  \includegraphics[width=0.88\textwidth]{fig7_ablation.pdf}
  \caption{Component ablation across 12 configurations. Left:
    $\Delta$\slo\% vs.\ Full RL; right: $\Delta$\etok\%.
    Red bars = $p<0.05$; gray bars = within noise; light gray
    band = $\pm 10\%$. Only continuous batching and pretraining
    show statistically significant effect; all individual AIGC
    reward and state components are within noise.}
  \label{fig:ablation}
\end{figure*}

\textbf{Reconciliation: how can RL Pareto-dominate baselines
yet be ablation-insensitive?} The apparent paradox between
\S\ref{sec:main-results} (RL dominates 10 baselines on means)
and Table~\ref{tab:ablation} (no single AIGC component is
significant) resolves as follows. Each baseline differs from
our scheduler in \emph{several} fundamental ways at once (no
learning; fixed cost functions without AIGC physics; a
different encoder or training recipe)---and falls behind by
7--28\%. When we ablate any \emph{single} AIGC component
within our scheduler, PPO compensates by extracting
equivalent signal from the remaining features. This tells us
\textbf{AIGC-aware scheduling is a holistic system
contribution, not a clever reward or state component}.

\subsection{Reward-Weight Sensitivity}
\label{sec:weight-sens}

The reward of Eq.~\ref{eq:reward} collapses a multi-objective
problem into a scalar weighted sum---are our conclusions an
artifact of one particular weight vector? We retrain from
scratch under four alternative configurations spanning the
design space---\emph{uniform} (all weights $1/7$),
\emph{time-heavy} (time term over half the mass),
\emph{AIGC-heavy}
(warm{+}batch{+}affinity raised to 0.70), and
\emph{base-heavy} (AIGC terms nearly removed)---each with
$N{=}30$ runs at the canonical point
(Table~\ref{tab:weight-sens}).

\begin{table}[t]
  \centering
  \caption{Reward-weight sensitivity (edge${=}5$,
    $\lambda{=}2$~req/s, $N{=}30$ per configuration). Weights
    listed in the order
    time/balance/match/warm/batch/affinity/cloud of
    Eq.~\ref{eq:reward}. $p$: two-sided Mann-Whitney U vs.\
    canonical, \slo{} / \etok. Every configuration
    Pareto-dominates the strongest baseline (PSO:
    \slo{}~0.235, \etok{}~2.81).}
  \label{tab:weight-sens}
  \scriptsize
  \setlength{\tabcolsep}{3pt}
  \begin{tabular}{@{}llccc@{}}
    \toprule
    Config & Weights & \slo{} & \etok{} & $p$ \\
    \midrule
    canonical & .25/.10/.10/.15/.20/.10/.10
      & $0.302 {\pm} 0.13$ & $2.61 {\pm} 0.67$ & --- \\
    uniform & $1/7$ each
      & $0.287 {\pm} 0.13$ & $2.61 {\pm} 0.68$ & .72 / .98 \\
    time-heavy & .55/.05/.05/.10/.10/.05/.10
      & $0.304 {\pm} 0.14$ & $2.60 {\pm} 0.67$ & .97 / .98 \\
    AIGC-heavy & .10/.05/.05/.25/.30/.15/.10
      & $0.300 {\pm} 0.12$ & $2.63 {\pm} 0.67$ & .86 / .96 \\
    base-heavy & .40/.20/.20/.05/.05/.05/.05
      & $0.305 {\pm} 0.14$ & $2.56 {\pm} 0.68$ & .89 / .74 \\
    \bottomrule
  \end{tabular}
\end{table}

Two observations. First, \textbf{the headline conclusion is
robust to the weight choice}: every configuration, including
the extreme ones, Pareto-dominates the strongest baseline
(\slo{} 0.287--0.305 vs.\ PSO's 0.235; \etok{} 2.56--2.63
vs.\ 2.81). Second, \textbf{no configuration differs
significantly from the canonical weights} (all $p \ge 0.72$):
even shifting the AIGC group's mass from 0.15 (base-heavy) to
0.70 (AIGC-heavy) leaves both objectives statistically
unchanged. This extends the ablation's null result from
\emph{removing} components to \emph{re-weighting} them, and
means our scalarization is benign: within a wide region of
weight space the learned policy lands on the same Pareto
point, so conclusions drawn from the canonical configuration
generalize.
